% BEDIENVORSCHAU FORMALER MODUS
% Öffentlicher Arbeitsablauf: eine Vorgabe, öffentliche Inhaltsbefehle und LetterBody.

\UseLetterForm{A}
\UsePreset{legal-de}

\UseRecipientManual
\SetManualRecipient{Fiktive Dokumentenstelle\\Allgemeine Korrespondenz\\Unterlagenweg 3\\12345 Beispielstadt}
\SetLetterSubject{Beispiel für Form A im Modus: legal}
\SetLetterOpening{Anrede}
\SetLetterClosing{Grußformel}

\SetRefIhrZeichen{DOK-12}
\SetRefMeinZeichen{MUSTER-26}
\SetRefIhreNachrichtVom{12. Mai 2026}
\SetLegalGeschaeftsnummer{}
\SetLegalSachverhalt{Musterunterlagen vom 12. Mai 2026}
\SetAttachmentsCustom{Übersicht der Musterunterlagen\\Fiktives Informationsblatt}

\begin{LetterBody}

\textbf{Was dieser Guide zeigt}

Für gegliederte formale Korrespondenz verwendet \texttt{legal-de} eine zurückhaltende schwarze Gestaltung und einen waagerechten Bezugsbereich. Diese Vorschau füllt nur den optionalen Sachverhaltsblock; der leere Geschäftsnummernblock entfällt. Der benutzerdefinierte Anlagenbefehl erzeugt die Liste am Ende, der Betreff bleibt direkt unter dem Bezugsbereich.

Passe Bezugswerte und Anlagen nur dort an, wo die Korrespondenz sie benötigt; der eigentliche Text bleibt in \texttt{LetterBody}. Die neutralisierte Vorschau beschreibt keinen realen Sachverhalt. Dies ist ein an DIN 5008 Form A orientiertes Layout für formale Korrespondenz, keine Rechtsberatung und nicht DIN-zertifiziert.

\end{LetterBody}
